\clearpage
\begin{explanationbox}
\textbf{\textcolor{CExplanation}{一、为什么会想到“半个周期”？——先看一个能画出来的函数}}
\begin{center}\begin{tikzpicture}[font=\small,>=stealth,line join=round]
\draw[gray!40,rounded corners=3pt] (0,0) rectangle (14.2,8.3);\node at (7.1,7.700000000000001) {周期2的余弦：轴每隔1出现，中心也每隔1出现};
\begin{scope}[shift={(2.3,3.8222222222222224)},x=1.2cm,y=1.277777777777778cm]\draw[->,gray] (-1,0)--(4,0) node[right]{$x$};\draw[->,gray] (0,-1.7)--(0,1.9) node[above]{$y$};\draw[blue,thick,solid,domain=-0.7:3.7,samples=180] plot(\x,{cos(deg(pi*\x))});\fill[orange!85!black] (0.5,0) circle(2.2pt);\fill[orange!85!black] (1.5,0) circle(2.2pt);\fill[orange!85!black] (2.5,0) circle(2.2pt);\fill[orange!85!black] (3.5,0) circle(2.2pt);\foreach \a in {0,1,2,3}{\draw[violet,dashed] (\a,-1.3)--(\a,1.3);\node[above] at (\a,1.35){$\a$};}\end{scope}\node[anchor=west,align=left,text width=5.0cm] at (8.4,4.0) {$f=\cos\pi x,\ T_0=2$\\[9pt]紫线：轴 $x=k$\\[9pt]橙点：中心 $(\frac12+k,0)$\\[9pt]同类间隔 $\frac{T_0}2=1$\\[9pt]但 $f(0)=1\ne f(1)=-1$};
\node[text width=12.6cm,align=center] at (7.1,.65) {半周期是同类对称位置的间隔，不表示每个函数值都在半周期后重复。};
\end{tikzpicture}\end{center}



取
\[
f(x)=\cos(\pi x).
\]
它每隔 $2$ 又回到原来的样子，最小正周期是 $2$。把图像画出来：



观察：
\begin{itemize}[leftmargin=*,itemsep=0.35em]
\item 图像从 $x=0$ 到 $x=2$，完成一整轮重复；
\item 但不用走完整个 $2$，只走一半到 $x=1$，又看到一条对称轴；
\item 位于 $x=\tfrac12,\tfrac32$ 的红点，则是图像绕其转半圈能重合的点；
\item 从对称轴 $x=0$ 走到最近的对称中心 $x=\tfrac12$，恰好是四分之一个最小正周期。
\end{itemize}

\textbf{特别注意：} $x=0$ 和 $x=1$ 虽然都是对称轴，但 $1$ \emph{不是} $\cos(\pi x)$ 的周期，因为 $f(0)=1$，$f(1)=-1$。对称轴每隔半周期出现，\textbf{不等于图像每隔半周期完全相同}。

\medskip

\end{explanationbox}
\clearpage
\begin{explanationbox}
\textbf{\textcolor{CExplanation}{二、用生活中的“折纸花纹”来理解}}
\begin{center}\begin{tikzpicture}[font=\small,>=stealth,line join=round]
\draw[gray!40,rounded corners=3pt] (0,0) rectangle (14.2,8.3);\node at (7.1,7.700000000000001) {证明一：右点先移回一个周期，再沿旧轴配对};
\begin{scope}[shift={(2.3857142857142857,3.8222222222222224)},x=0.8571428571428571cm,y=1.277777777777778cm]\draw[->,gray] (-1.5,0)--(5.5,0) node[right]{$x$};\draw[->,gray] (0,-1.7)--(0,1.9) node[above]{$y$};\draw[blue,thick,solid,domain=-1:5,samples=180] plot(\x,{cos(deg(pi*(\x-1)/2))});\fill[orange!85!black] (-0.5,-0.7071067811865476) circle(2.2pt);\fill[orange!85!black] (2.5,-0.7071067811865476) circle(2.2pt);\fill[orange!85!black] (3.5,-0.7071067811865476) circle(2.2pt);\foreach \a in {1,3}{\draw[violet,dashed] (\a,-1.25)--(\a,1.25);\node[above] at (\a,1.3){$x=\a$};}\foreach \X/\L in {-.5/{R-T},2.5/L,3.5/R}{\draw[gray,dashed] (\X,-.70710678)--(\X,-1.2);\node[below=4pt] at (\X,-1.4){$\L$};}\end{scope}\node[anchor=west,align=left,text width=5.0cm] at (8.4,4.0) {$T=4,\ a=1,\ t=\frac12$\\[9pt]新轴 $a+T/2=3$\\[9pt]$R=3.5,\ R-T=-0.5$\\[9pt]$L=2.5$\\[9pt]三处都为 $-\frac{\sqrt2}2$};
\node[text width=12.6cm,align=center] at (7.1,.65) {图取一组具体距离；证明中t任意，平移一次与沿旧轴配对就得到新轴。};
\end{tikzpicture}\end{center}



想象走廊墙上有一串不断重复的花纹，每段长 $T$；又知道其中一条竖直线是一条“折痕”，花纹沿着它对折可以重合。

在已知折痕左边取一个位置。沿着“对折”和“向右移动一个完整花纹”这两条规则，可以在折痕右边找到相同的花纹。因为“整段重复”不分从哪里开始，原折痕向右走半段 $T/2$ 的地方同样能对折重合。再重复，就得到一串折痕。

对称中心可以想成\textbf{把纸绕钉在某一点的图钉旋转 $180^\circ$}，原来处于中心上方的点会落到下方，中心高度始终不变。这就是为什么新中心的横坐标改变，纵坐标 $c$ 不变。

\medskip
\end{explanationbox}\clearpage\begin{explanationbox}\textbf{\textcolor{CExplanation}{再看一个初中生也能画的“折线波”}}
\begin{center}\begin{tikzpicture}[font=\small,>=stealth,line join=round]
\draw[gray!40,rounded corners=3pt] (0,0) rectangle (14.2,8.3);\node at (7.1,7.700000000000001) {不用三角函数：折线波也有半周期轴和四分之一周期中心};
\begin{scope}[shift={(3.242857142857143,2.4863636363636363)},x=0.8571428571428571cm,y=2.0909090909090904cm]\draw[->,gray] (-2.5,0)--(4.5,0) node[right]{$x$};\draw[->,gray] (0,-0.4)--(0,1.8) node[above]{$y$};\draw[blue,thick,solid,domain=-2:-1,samples=180] plot(\x,{1*\x+(2)});\draw[blue,thick,solid,domain=-1:0,samples=180] plot(\x,{-1*\x+(0)});\draw[blue,thick,solid,domain=0:1,samples=180] plot(\x,{1*\x+(0)});\draw[blue,thick,solid,domain=1:2,samples=180] plot(\x,{-1*\x+(2)});\draw[blue,thick,solid,domain=2:3,samples=180] plot(\x,{1*\x+(-2)});\draw[blue,thick,solid,domain=3:4,samples=180] plot(\x,{-1*\x+(4)});\fill[orange!85!black] (0.5,0.5) circle(2.2pt);\fill[orange!85!black] (1.5,0.5) circle(2.2pt);\fill[orange!85!black] (2.5,0.5) circle(2.2pt);\fill[orange!85!black] (3.5,0.5) circle(2.2pt);\foreach \a in {-2,-1,0,1,2,3,4}{\draw[violet,dashed] (\a,0)--(\a,1.2);}\end{scope}\node[anchor=west,align=left,text width=5.0cm] at (8.4,4.0) {周期 $T_0=2$\\[9pt]轴 $x=0,1,2,\ldots$\\[9pt]中心 $(\frac12+k,\frac12)$\\[9pt]所有中心高度都为 $\frac12$\\[9pt]原理与余弦图相同};
\node[text width=12.6cm,align=center] at (7.1,.65) {实际由一段上升直线和一段下降直线周期延拓；两端值均为0，接头相容。};
\end{tikzpicture}\end{center}



不用正弦、余弦，只用直线：先从 $(0,0)$ 走到 $(1,1)$，再走到 $(2,0)$；每隔 $2$ 个单位把这段折线照样重复。用式子写就是
\[
f(x)=\begin{cases}
x,&0\le x\le1,\\
2-x,&1<x\le2,
\end{cases}\qquad f(x+2)=f(x).
\]



这里的周期是 $2$，对称轴在 $x=0,1,2,\ldots$，相邻两轴间距就是 $1=2/2$。红点如 $(\tfrac12,\tfrac12)$ 是对称中心，离最近的轴只有 $\tfrac12=2/4$。直线折成的“山峰、山谷”，同样能出现\textbf{半周期对称轴、四分之一周期对称中心}，所以这不是三角函数独有的巧合。

\medskip

\end{explanationbox}
\clearpage
\begin{explanationbox}
\textbf{\textcolor{CExplanation}{三、最简单的数字演示：不用三角函数也能看懂}}
\begin{center}\begin{tikzpicture}[font=\small,>=stealth,line join=round]
\draw[gray!40,rounded corners=3pt] (0,0) rectangle (14.2,8.3);\node at (7.1,7.700000000000001) {证明一：右点先移回一个周期，再沿旧轴配对};
\begin{scope}[shift={(2.3857142857142857,3.8222222222222224)},x=0.8571428571428571cm,y=1.277777777777778cm]\draw[->,gray] (-1.5,0)--(5.5,0) node[right]{$x$};\draw[->,gray] (0,-1.7)--(0,1.9) node[above]{$y$};\draw[blue,thick,solid,domain=-1:5,samples=180] plot(\x,{cos(deg(pi*(\x-1)/2))});\fill[orange!85!black] (-0.5,-0.7071067811865476) circle(2.2pt);\fill[orange!85!black] (2.5,-0.7071067811865476) circle(2.2pt);\fill[orange!85!black] (3.5,-0.7071067811865476) circle(2.2pt);\foreach \a in {1,3}{\draw[violet,dashed] (\a,-1.25)--(\a,1.25);\node[above] at (\a,1.3){$x=\a$};}\foreach \X/\L in {-.5/{R-T},2.5/L,3.5/R}{\draw[gray,dashed] (\X,-.70710678)--(\X,-1.2);\node[below=4pt] at (\X,-1.4){$\L$};}\end{scope}\node[anchor=west,align=left,text width=5.0cm] at (8.4,4.0) {$T=4,\ a=1,\ t=\frac12$\\[9pt]新轴 $a+T/2=3$\\[9pt]$R=3.5,\ R-T=-0.5$\\[9pt]$L=2.5$\\[9pt]三处都为 $-\frac{\sqrt2}2$};
\node[text width=12.6cm,align=center] at (7.1,.65) {图取一组具体距离；证明中t任意，平移一次与沿旧轴配对就得到新轴。};
\end{tikzpicture}\end{center}



假设一个函数满足
\[
f(x+4)=f(x),\qquad f(1+t)=f(1-t).
\]
已知 $x=1$ 是对称轴，问 $x=3$ 是不是？从 $3$ 两边随意取距离 $0.5$ 的两个点：
\[
\begin{aligned}
f(3.5)&=f(-0.5) &&\text{（向左移动 $4$）}\\
       &=f(2.5) &&\text{（$-0.5$ 与 $2.5$ 关于 $1$ 对称）}.
\end{aligned}
\]
所以 $3.5$ 和 $2.5$ 的函数值相同。换成任意距离 $t$ 也同样成立，因此 $x=3$ 确实是对称轴。

\textbf{这里最关键的理解：}题目中的两个条件都是“\emph{对任意} $x$ 成立”，所以才能把 $0.5$ 换成任意距离。只核对一对点，不能证明整条直线是对称轴。

\medskip

\end{explanationbox}
\clearpage
\begin{explanationbox}
\textbf{\textcolor{CExplanation}{四、为什么“两次对折 = 向一边移动两倍距离”？}}
\begin{center}\begin{tikzpicture}[font=\small,>=stealth,line join=round]
\draw[gray!40,rounded corners=3pt] (0,0) rectangle (14.2,8.3);\node at (7.1,7.700000000000001) {例4：关于负1与3反射两次，横向总共移动8};
\begin{scope}[shift={(2.7153846153846155,3.8222222222222224)},x=0.46153846153846156cm,y=1.277777777777778cm]\draw[->,gray] (-3.5,0)--(9.5,0) node[right]{$x$};\draw[->,gray] (0,-1.7)--(0,1.9) node[above]{$y$};\draw[blue,thick,solid,domain=-3:9,samples=180] plot(\x,{cos(deg(pi*(\x+1)/4))});\fill[orange!85!black] (0,0.7071067811865476) circle(2.2pt);\fill[orange!85!black] (-2,0.7071067811865476) circle(2.2pt);\fill[orange!85!black] (8,0.7071067811865476) circle(2.2pt);\foreach \a in {-1,3}{\draw[violet,dashed] (\a,-1.3)--(\a,1.3);\node[above] at (\a,1.35){$x=\a$};}\end{scope}\node[anchor=west,align=left,text width=5.0cm] at (8.4,4.0) {$f=\cos\frac{\pi(x+1)}4$\\[9pt]轴 $x=-1,3$\\[9pt]$0\to-2\to8$\\[9pt]三点高度 $\frac{\sqrt2}2$\\[9pt]周期可取 $2\times4=8$};
\node[text width=12.6cm,align=center] at (7.1,.65) {两次反射的合成位移，是两轴间距的两倍；这里只证明一个周期。};
\end{tikzpicture}\end{center}



在数轴上设两条折痕分别在 $a_1$、$a_2$，任取一个点 $u$：
\[
\underbrace{u}_{\text{起点}}
\ \xrightarrow{\text{沿 $a_1$ 对折}}\
2a_1-u
\ \xrightarrow{\text{再沿 $a_2$ 对折}}\
2a_2-(2a_1-u)=u+2(a_2-a_1).
\]

例如折痕在 $x=1$ 与 $x=3$，点 $u=0$：第一次变成 $2$，第二次变成 $4$。\textbf{两次对折后位置总共向右移动了 $4$，正好是折痕间距 $2$ 的两倍。} 如果函数图像两次对折后都不变，那么向右移动 $4$ 也必定不变，所以 $4$ 是一个周期。

\medskip

\end{explanationbox}
\clearpage
\begin{explanationbox}
\textbf{\textcolor{CExplanation}{五、眼前一亮：为什么“对称轴到对称中心”常常是四分之一周期？}}
\begin{center}\begin{tikzpicture}[font=\small,>=stealth,line join=round]
\draw[gray!40,rounded corners=3pt] (0,0) rectangle (14.2,8.3);\node at (7.1,7.700000000000001) {轴到中心不一定最近：也可能跨3个四分之一周期};
\begin{scope}[shift={(2.158823529411765,3.8222222222222224)},x=1.7647058823529411cm,y=1.277777777777778cm]\draw[->,gray] (-0.6,0)--(2.8,0) node[right]{$x$};\draw[->,gray] (0,-1.7)--(0,1.9) node[above]{$y$};\draw[blue,thick,solid,domain=-0.3:2.5,samples=180] plot(\x,{cos(deg(pi*\x))});\fill[orange!85!black] (0,1) circle(2.2pt);\fill[orange!85!black] (0.5,0) circle(2.2pt);\fill[orange!85!black] (1,-1) circle(2.2pt);\fill[orange!85!black] (1.5,0) circle(2.2pt);\draw[violet,dashed] (0,-1.3)--(0,1.3);\foreach \v/\L in {.5/\frac12,1/1,1.5/\frac32,2/2}{\node[below=4pt] at (\v,-1.35){$\L$};}\end{scope}\node[anchor=west,align=left,text width=5.0cm] at (8.4,4.0) {$T_0=2$，轴 $x=0$\\[9pt]中心 $(\frac12,0),(\frac32,0)$\\[9pt]$\frac12=T_0/4$\\[9pt]$\frac32=3T_0/4$\\[9pt]一般为奇数个 $T_0/4$};
\node[text width=12.6cm,align=center] at (7.1,.65) {四分之一周期关系使用最小正周期；不能把较远中心也误当最近的中心。};
\end{tikzpicture}\end{center}



仍看 $\cos(\pi x)$：在 $x=0$ 处是波峰，在 $x=\tfrac12$ 处恰好经过中间高度 $y=0$，在 $x=1$ 处是波谷。
\[
\underbrace{0}_{\text{对称轴（峰）}}
\ \longrightarrow\
\underbrace{\tfrac12}_{\text{对称中心（中线）}}
\ \longrightarrow\
\underbrace{1}_{\text{对称轴（谷）}}
\]
完整一轮长 $2$，所以峰到中线长 $\tfrac12=2/4$。一般函数也是这样，但要有\textbf{一条对称轴和一个对称中心}，并且要用\textbf{最小正周期}才能描述“奇数个四分之一周期”；详细证明放在 03。

\medskip

\end{explanationbox}
\clearpage
\begin{explanationbox}
\textbf{\textcolor{CExplanation}{六、把本结论和原来学过的知识连起来}}
\begin{center}\begin{tikzpicture}[font=\small,>=stealth,line join=round]
\draw[gray!40,rounded corners=3pt] (0,0) rectangle (14.2,8.3);\node at (7.1,7.700000000000001) {周期2的余弦：轴每隔1出现，中心也每隔1出现};
\begin{scope}[shift={(2.3,3.8222222222222224)},x=1.2cm,y=1.277777777777778cm]\draw[->,gray] (-1,0)--(4,0) node[right]{$x$};\draw[->,gray] (0,-1.7)--(0,1.9) node[above]{$y$};\draw[blue,thick,solid,domain=-0.7:3.7,samples=180] plot(\x,{cos(deg(pi*\x))});\fill[orange!85!black] (0.5,0) circle(2.2pt);\fill[orange!85!black] (1.5,0) circle(2.2pt);\fill[orange!85!black] (2.5,0) circle(2.2pt);\fill[orange!85!black] (3.5,0) circle(2.2pt);\foreach \a in {0,1,2,3}{\draw[violet,dashed] (\a,-1.3)--(\a,1.3);\node[above] at (\a,1.35){$\a$};}\end{scope}\node[anchor=west,align=left,text width=5.0cm] at (8.4,4.0) {$f=\cos\pi x,\ T_0=2$\\[9pt]紫线：轴 $x=k$\\[9pt]橙点：中心 $(\frac12+k,0)$\\[9pt]同类间隔 $\frac{T_0}2=1$\\[9pt]但 $f(0)=1\ne f(1)=-1$};
\node[text width=12.6cm,align=center] at (7.1,.65) {半周期是同类对称位置的间隔，不表示每个函数值都在半周期后重复。};
\end{tikzpicture}\end{center}


\begin{itemize}[leftmargin=*,itemsep=0.4em]
\item \textbf{与偶函数相连：} $f(-x)=f(x)$ 就是对称轴在 $x=0$。已知一个周期 $T$，立刻知道 $x=kT/2$ 都是对称轴。
\item \textbf{与奇函数相连：} $f(-x)=-f(x)$ 就是对称中心在 $(0,0)$。已知一个周期 $T$，立刻知道 $(kT/2,0)$ 都是对称中心。
\item \textbf{与函数平移相连：} $f(a+t)=f(a-t)$ 其实是把对称轴移到 $x=a$；而 $f(b+t)+f(b-t)=2c$ 是把中心移到 $(b,c)$。
\item \textbf{与抽象函数求值相连：} 不必知道 $f$ 的公式。等式 $f(x+T)=f(x)$ 负责“搬位置”，对称关系负责“换方向”。把难算的自变量搬回已知值附近，就能求解。
\end{itemize}
\end{explanationbox}
